\documentclass[twocolumn,preprintnumbers,amsmath,amssymb,prd]{revtex4-2}

\usepackage{graphicx}
\usepackage{amsmath,amssymb}
\usepackage{bm}
\usepackage{hyperref}
\usepackage{xcolor}
\usepackage{booktabs}

\begin{document}

\title{Cosmological Predictions from Topological Phase Transitions:\\Vacuum Energy, CMB Temperature, and Neutrino Mass\\from $A_2$ BKT Projections}

\author{Jesse C.P. Won Ting}
\email{jass168611@gmail.com}
\affiliation{Independent Researcher}

\date{\today}

\begin{abstract}
We demonstrate that four fundamental physical quantities spanning 96 orders of magnitude---the dark matter relaxation time ($\sim 10^{41}\,\tau_0$), the vacuum energy density ($\rho_\Lambda/\Lambda^4 \sim 10^{-55}$), the CMB temperature ($T_{\rm CMB}/\Lambda \sim 10^{-14}$), and the neutrino-to-electron mass ratio ($m_\nu/m_e \sim 10^{-7}$)---all derive from a single $2\times2$ Cartan matrix $K_{A_2} = \bigl(\begin{smallmatrix}2 & -1\\ -1 & 2\end{smallmatrix}\bigr)$ with eigenvalues $\lambda_1 = 1$ and $\lambda_2 = 3$, combined with the BKT reduced temperature $t = 1.1\times10^{-3}$ determined by atomic spectroscopy.
The four quantities differ only in the \emph{projection rule} applied to the two BKT parameters $b_1 = \pi/2$ and $b_2 = \pi/6$: serial relaxation (dark matter), fugacity product (vacuum energy), equal-weight average (CMB), and single-mode decoupling (neutrino mass).
A uniqueness theorem proves that $A_2$ ($\mathrm{SU}(3)$) is the only non-trivial rank-2 root system for which the BKT parameter sum equals the root angle, simultaneously explaining why there are exactly three colors and why the vacuum is stable.
We further resolve the apparent tension between $b = \pi$ (dark matter relaxation) and $b = 2\pi/3$ (vacuum energy) as projections of the same Cartan matrix in different bases: the Cartesian (field) basis gives $b_{\rm relax} = \pi$; the eigenbasis gives $b_{\rm vac} = 2\pi/3$.
The neutrino mass $m_\nu = m_e \exp(-\pi/(6\sqrt{t})) = 0.072\,$eV is a zero-parameter prediction, consistent with oscillation data ($\sim 0.05\,$eV) and the Planck bound ($\Sigma m_\nu < 0.12\,$eV).
Applied to baryogenesis, the BKT critical slowing down at $t_{\rm freeze} = 1.33\times10^{-3} \approx t_{\rm ESFT}$ satisfies all three Sakharov conditions, yielding $\eta_B \sim 10^{-8}$ (within 1.5 orders of magnitude of observation).
No adjustable parameters appear at any step.
\end{abstract}

\maketitle

%% ============================================================
\section{Introduction}
\label{sec:intro}
%% ============================================================

The cosmological constant problem ($\rho_\Lambda/M_{\rm Pl}^4 \sim 10^{-120}$), the neutrino mass hierarchy ($m_\nu/m_e \sim 10^{-6}$), the CMB temperature ($T_{\rm CMB} \sim 3\,$K in a universe with characteristic energy scale $\sim 10^{14}\,$K), and the dark matter relaxation scale ($\tau_{\rm DM}/\tau_{\rm micro} \sim 10^{41}$) are conventionally treated as independent puzzles requiring separate solutions.

In this paper, we show that all four quantities arise from a single mathematical object: the $A_2$ Cartan matrix $K = \bigl(\begin{smallmatrix}2 & -1\\ -1 & 2\end{smallmatrix}\bigr)$, evaluated in the BKT-ordered phase at $t = 1.1\times10^{-3}$.
The only difference between the four predictions is the \emph{projection rule}---how the two eigenmode BKT parameters $b_1$ and $b_2$ are combined.
This structure is analogous to Bell's theorem, where a single quantum state yields different correlations depending on the measurement basis.

The paper is organized as follows.
Section~\ref{sec:cartan} establishes the $A_2$ eigenmode decomposition.
Section~\ref{sec:uniqueness} proves the uniqueness theorem.
Section~\ref{sec:four_projections} derives all four quantities from their respective projection rules.
Section~\ref{sec:b_resolution} resolves the $b = \pi$ vs.\ $2\pi/3$ tension.
Section~\ref{sec:neutrino} derives the neutrino mass.
Section~\ref{sec:baryogenesis} presents the baryogenesis estimate.
Sections~\ref{sec:hubble}--\ref{sec:strongCP} present the Hubble constant, dark matter, and strong CP results.
Section~\ref{sec:bell} discusses the Bell analogy.
Section~\ref{sec:tests} summarizes testable predictions.

%% ============================================================
\section{$A_2$ Eigenmode Decomposition}
\label{sec:cartan}
%% ============================================================

The three ESFT phase fields $\Theta_a$ ($a = 1,2,3$) satisfy the $A_2$ constraint $\Theta_1 + \Theta_2 + \Theta_3 = 0$, leaving two independent fields.
The $A_2$ Cartan matrix
\begin{equation}
K = \begin{pmatrix} 2 & -1 \\ -1 & 2 \end{pmatrix}
\label{eq:cartan}
\end{equation}
governs the coupling between them, with eigenvalues $\lambda_1 = 1$ (soft mode) and $\lambda_2 = 3$ (stiff mode).

\subsection{Eigenmodes}

The eigenvectors are $\Theta_\pm = (\Theta_1 \pm \Theta_2)/\sqrt{2}$, and the Lagrangian diagonalizes as:
\begin{equation}
\mathcal{L} = \sum_{i=1}^{2}\left[\frac{f^2\lambda_i}{2}(\partial\phi_i)^2 - V(\phi_i)\right],
\end{equation}
where $\phi_i$ are the eigenmodes.
Each eigenmode experiences an effective BKT stiffness $K_{\rm eff,i} = \beta f^2 \lambda_i$, modifying its BKT parameter.

\subsection{BKT parameters}

Each eigenmode undergoes an independent BKT transition with parameter:
\begin{equation}
b_i = \frac{\pi}{2\lambda_i}\,.
\label{eq:bi}
\end{equation}
Explicitly: $b_1 = \pi/2$ (standard BKT for $\lambda_1 = 1$) and $b_2 = \pi/6$ (for $\lambda_2 = 3$).
The derivation of $b_2$ is presented in Section~\ref{sec:uniqueness}.

%% ============================================================
\section{Uniqueness Theorem: $A_2$ Selects Three Colors}
\label{sec:uniqueness}
%% ============================================================

\textbf{Theorem.} \emph{Among all finite-type rank-2 root systems, $A_2$ ($c=1$) is the unique non-trivial solution of $b_1 + b_2 = \theta_{\rm root}$, where $\theta_{\rm root}$ is the angle between simple roots.}

For a general rank-2 Cartan matrix with off-diagonal element $-c$, the eigenvalues are $\lambda_{1,2} = 2 \mp c$, giving:
\begin{equation}
b_1 + b_2 = \frac{\pi}{2}\cdot\frac{\mathrm{Tr}(K)}{\det(K)} = \frac{2\pi}{4-c^2}\,.
\label{eq:bsum}
\end{equation}
The root angle is $\theta_{\rm root} = \arccos(-c/2)$.

\emph{Proof.}
\begin{itemize}
\item $c = 0$ ($A_1 \times A_1$): $b_{\rm sum} = \pi/2 = \theta_{\rm root}$ (trivial---two decoupled BKT systems).
\item $c = 1$ ($A_2 = \mathrm{SU}(3)$): $b_{\rm sum} = 2\pi/3 = \theta_{\rm root}$ (unique non-trivial solution).
\item $c \geq 2$: $4 - c^2 \leq 0 \Rightarrow b_{\rm sum}$ diverges; no stable BKT phase exists.
\end{itemize}
\hfill$\square$

\textbf{Corollary:} $b_2 = \pi/6$ follows from two premises alone: (1)~$b_1 = \pi/2$ (standard BKT), and (2)~$b_1 + b_2 = 2\pi/3$ (uniqueness theorem).
The functional form $b_i = \pi/(2\lambda_i)$ is determined by the two data points $(\lambda_1, b_1) = (1, \pi/2)$ and $(\lambda_2, b_2) = (3, \pi/6)$.

\textbf{Physical significance:} Three colors is not an arbitrary choice---it is the unique gauge structure compatible with vacuum stability under BKT ordering.

%% ============================================================
\section{Four Projections, 96 Orders of Magnitude}
\label{sec:four_projections}
%% ============================================================

All four quantities share the same BKT exponential structure $\exp(\pm b_{\rm eff}/\sqrt{t})$ with $t = 1.1\times10^{-3}$.
They differ only in the projection rule that determines $b_{\rm eff}$:

\begin{center}
\begin{tabular}{lllll}
\toprule
Physical quantity & Projection rule & $b_{\rm eff}$ & Exponent & Scale \\
\midrule
DM relaxation & Cartesian serial & $\pi$ & $10^{+41}$ & $\tau \sim 10^{15}\,$s \\
Vacuum energy & Eigenmode product & $4\pi/3$ & $10^{-55}$ & $\rho_\Lambda/\Lambda^4$ \\
CMB temperature & Eigenmode average & $\pi/3$ & $10^{-14}$ & $T_{\rm CMB}/\Lambda$ \\
Neutrino mass & Missing $\lambda_2$ mode & $\pi/6$ & $10^{-7}$ & $m_\nu/m_e$ \\
\bottomrule
\end{tabular}
\end{center}

\subsection{Why each projection rule}

\textbf{Dark matter relaxation ($b = \pi$):}
Each of the two independent fields relaxes to equilibrium through its own BKT dynamics.
In the Cartesian (field) basis, each field sees diagonal stiffness $K_{aa} = 2$, giving $b_{\rm single} = \pi/2$.
Serial relaxation (field 1 then field 2) multiplies the timescales: $b_{\rm relax} = 2 \times \pi/2 = \pi$. (Section~\ref{sec:b_resolution}.)

\textbf{Vacuum energy ($b = 4\pi/3$):}
The vacuum state requires \emph{both} eigenmodes to be simultaneously vortex-free.
Independent modes: fugacities multiply.
Each mode contributes $y_i^2 = \exp(-2b_i/\sqrt{t})$, so:
\begin{equation}
y_1^2 y_2^2 = \exp\!\left(-\frac{2(b_1+b_2)}{\sqrt{t}}\right) = \exp\!\left(-\frac{4\pi/3}{\sqrt{t}}\right).
\end{equation}

\textbf{CMB temperature ($b = \pi/3$):}
The electromagnetic field $\Theta_{\rm EM} = \Theta_1$ projects equally onto both eigenmodes (Section~\ref{sec:em_projection}).
The BKT suppression is the equal-weight average: $b_{\rm EM} = (b_1 + b_2)/2 = \pi/3$.

\textbf{Neutrino mass ($b = \pi/6$):}
Neutrinos are color singlets---they do not couple to the $\lambda_2 = 3$ (stiff) eigenmode.
Only the $\lambda_1 = 1$ mode provides BKT protection.
The missing mode contributes a suppression factor $\exp(-b_2/\sqrt{t}) = \exp(-\pi/(6\sqrt{t}))$ relative to charged leptons. (Section~\ref{sec:neutrino}.)

\subsection{The 96-order span}

From $\tau/\tau_0 \sim 10^{+41}$ (dark matter) to $\rho_\Lambda/\Lambda^4 \sim 10^{-55}$ (vacuum energy), the four quantities span:
\begin{equation}
10^{+41} \div 10^{-55} = 10^{96}\quad\text{orders of magnitude,}
\end{equation}
all determined by one $2\times 2$ matrix and one parameter $t$.

%% ============================================================
\section{Resolution of $b = \pi$ vs.\ $2\pi/3$}
\label{sec:b_resolution}
%% ============================================================

Previous ESFT papers used $b = \pi$ for dark matter relaxation (Paper~I) and $b_1 + b_2 = 2\pi/3$ for vacuum energy.
This apparent tension is now resolved: they are projections of the same Cartan matrix in different bases.

\subsection{Cartesian basis (field relaxation)}

Each field $\Theta_a$ relaxes locally, governed by its diagonal stiffness $K_{aa} = 2$.
The off-diagonal coupling $K_{ab} = -1$ determines the equilibrium configuration but not the relaxation rate.
Each field independently gives $b = \pi/2$ (standard BKT with $K = 2$).
Two independent fields in series: $b_{\rm relax} = 2 \times \pi/2 = \pi$.

\subsection{Eigenbasis (vortex fugacity)}

Vortices are collective excitations of the \emph{eigenmodes}, not individual fields.
Eigenmode $i$ has effective stiffness $\propto \lambda_i$, giving $b_i = \pi/(2\lambda_i)$.
The vacuum fugacity is a product over eigenmodes: $b_{\rm vac} = 2(b_1 + b_2) = 4\pi/3$, or per-mode sum $b_1 + b_2 = 2\pi/3$.

\subsection{Why different bases}

The physical distinction is:
\begin{itemize}
\item \textbf{Relaxation} = return to equilibrium from a perturbed state. Each field relaxes \emph{locally} in the Cartesian basis; inter-field coupling only shifts the equilibrium point.
\item \textbf{Vortex excitation} = topological fluctuation in the ordered phase. Vortices are eigenmodes of the collective stiffness matrix; their fugacity depends on $\lambda_i$.
\end{itemize}

Both readings are correct projections of the same matrix $K = (2,-1;-1,2)$:
diagonal elements $= 2$ (field relaxation), eigenvalues $= \{1,3\}$ (vortex fugacity).

%% ============================================================
\section{Electromagnetic Projection}
\label{sec:em_projection}
%% ============================================================

The electromagnetic phase is $\Theta_{\rm EM} = Q_u\Theta_1 + Q_d\Theta_2 + Q_s\Theta_3 = \Theta_1$ (after applying the $A_2$ constraint with $Q_u = 2/3$, $Q_d = Q_s = -1/3$).
Decomposing into eigenmodes: $\Theta_{\rm EM} = (\Theta_+ + \Theta_-)/\sqrt{2}$, with equal weight $1/2$ on each.

The BKT suppression of the EM sector is:
\begin{equation}
\boxed{b_{\rm EM} = \frac{b_1 + b_2}{2} = \frac{\pi/2 + \pi/6}{2} = \frac{\pi}{3}}\,.
\label{eq:b_em}
\end{equation}

%% ============================================================
\section{CMB Temperature}
\label{sec:cmb}
%% ============================================================

The soliton core defines $\Lambda_{\rm core} = \hbar c/\varkappa \approx 73\,$GeV.
The BKT suppression of the EM sector gives:
\begin{equation}
\boxed{T_{\rm CMB} = \frac{\Lambda_{\rm core}}{2\pi k_B} \times \exp\!\left(-\frac{\pi}{3\sqrt{t}}\right)}
\label{eq:Tcmb}
\end{equation}

With $\varkappa_{\rm phys} = 0.0027\,$fm:
\begin{equation}
T_{\rm CMB}^{\rm ESFT} = 2.62\,\text{K}\qquad(\text{obs.\ }2.725\,\text{K},\;\;3.9\%).
\end{equation}
All inputs are from ESFT Papers~I--II; no cosmological parameters are used.

%% ============================================================
\section{Vacuum Energy from BKT Fugacity}
\label{sec:vacuum}
%% ============================================================

The vacuum energy is the residual disorder energy from bound vortex-antivortex pairs:
\begin{equation}
\rho_{\rm vac} = \Lambda^4\,y_1^2\,y_2^2 = \Lambda^4\exp\!\left(-\frac{2(b_1+b_2)}{\sqrt{t}}\right).
\end{equation}

\textbf{Numerical result:}
\begin{align}
\frac{2(b_1+b_2)}{\sqrt{t}} &= \frac{2\times 2\pi/3}{0.03317} = \frac{4\pi/3}{0.03317} = 126.3\,, \notag\\
\rho_\Lambda^{\rm ESFT} &= (73\,\text{GeV})^4 \times 10^{-54.85} \notag\\
&= \boxed{4.0\times10^{-48}\,\text{GeV}^4}\,.
\label{eq:rho_lambda}
\end{align}
Observed: $\rho_\Lambda^{\rm obs} = 3.6\times10^{-47}\,$GeV$^4$~\cite{Planck2020}.
The BKT mechanism suppresses the bare value by 55 orders of magnitude, within a factor of 9 of observation.

The factor-9 residual corresponds to a $3.5\%$ shift in $t$ (from $1.1\times10^{-3}$ to $1.14\times10^{-3}$).
This sensitivity is \emph{the key testable prediction}: an independent determination of $t$ to 1\% precision via precision spectroscopy will pin down $\rho_\Lambda$ within a factor of 2.

%% ============================================================
\section{Neutrino Mass from BKT Color Suppression}
\label{sec:neutrino}
%% ============================================================

\subsection{Physical argument}

Charged leptons couple to the full $A_2$ structure through electromagnetic interactions.
Neutrinos are color singlets---they inhabit the $\mathrm{SU}(2)_{\rm weak}$ subspace (rank 1 of $A_2$'s rank 2).
The $\lambda_2 = 3$ eigenmode is the ``pure color'' BKT mode; neutrinos do not participate in its topological ordering.

This missing BKT protection introduces an additional suppression factor:
\begin{equation}
\frac{m_\nu}{m_e} = \exp\!\left(-\frac{b_2}{\sqrt{t}}\right) = \exp\!\left(-\frac{\pi}{6\sqrt{t}}\right).
\end{equation}

\subsection{Derivation}

The electron mass in ESFT is $m_e = \sigma/(2\pi\Lambda) = 0.511\,$MeV (Paper~III, 0.24\% accuracy).
The neutrino mass:
\begin{equation}
\boxed{m_\nu = m_e \times \exp\!\left(-\frac{\pi}{6\sqrt{t}}\right) = 0.072\,\text{eV}}
\label{eq:m_nu}
\end{equation}
using $\pi/(6\sqrt{t}) = \pi/(6\times0.03317) = 15.77$.

\subsection{Comparison with experiment}

\begin{center}
\begin{tabular}{lccc}
\toprule
Quantity & ESFT & Experiment & Note \\
\midrule
$m_{\nu_3}$ & 0.072\,eV & $\sim 0.050$\,eV & 44\%, tree-level \\
$\Sigma m_\nu$ & $\sim 0.08$\,eV & $< 0.12$\,eV & Planck bound~\cite{Planck2020} \\
Hierarchy & Normal & TBD & JUNO ($\sim$2027) \\
$\Delta m^2_{32}$ & $\sim 0.003\,$eV$^2$ & 0.00245\,eV$^2$ & $\sim$20\% \\
\bottomrule
\end{tabular}
\end{center}

The 44\% deviation is consistent with tree-level accuracy: the exponential sensitivity to $t$ means that a $+9\%$ shift in $t$ (from $1.1\times10^{-3}$ to $1.2\times10^{-3}$) shifts $m_\nu$ from 0.072 to 0.048\,eV, matching observation exactly.
This is the \emph{same direction} as the vacuum energy $t$-shift ($+3.5\%$), suggesting $t_{\rm phys}$ is slightly above $1.1\times10^{-3}$.

\subsection{Three-generation splitting}

The large splitting ($\Delta m^2_{32} \sim 10^{-3}\,$eV$^2$) arises from the first Fibonacci step of the weak ladder (Paper~VII).
The small splitting ($\Delta m^2_{21} \sim 10^{-5}\,$eV$^2$) is a one-loop effect.
The ratio $\Delta m^2_{21}/\Delta m^2_{32} \sim 0.03$ reflects the tree/loop hierarchy $\sim \alpha_W/(4\pi) \sim 10^{-2}$.

\subsection{Forward prediction}

ESFT predicts \textbf{normal hierarchy} ($m_3 > m_2 > m_1$), testable by JUNO (expected $\sim$2027).
If inverted hierarchy is observed, this prediction is falsified.

%% ============================================================
\section{Hubble Constant}
\label{sec:hubble}
%% ============================================================

The Hubble parameter:
\begin{equation}
H_0 = \frac{c\,\varepsilon}{N_c\,\varkappa\,e^{\pi/\sqrt{t}}} \times \sqrt{\frac{\pi}{2}} \times \sqrt{1+\frac{2}{N_c^2}} = 72.5\,\text{km/s/Mpc}\,,
\label{eq:H0}
\end{equation}
between Planck ($67.4 \pm 0.5$)~\cite{Planck2020} and SH0ES ($73.0 \pm 1.0$)~\cite{Riess2022}.
The $\varkappa$-dependence cancels identically: $H_0 = \sqrt{\sigma}/e^{\pi/\sqrt{t}} \times$ factors---the most robust ESFT prediction.

%% ============================================================
\section{Strong CP Problem}
\label{sec:strongCP}
%% ============================================================

The QCD instanton amplitude is modulated by the same two-mode BKT fugacity:
\begin{equation}
\theta_{\rm eff} = \theta_{\rm bare} \times \exp\!\left(-\frac{2(b_1+b_2)}{\sqrt{t}}\right) \sim \theta_{\rm bare}\times 10^{-55}\,.
\end{equation}
Even with $\theta_{\rm bare} = O(1)$: $|\theta_{\rm eff}| \sim 10^{-55} \ll 10^{-10}$, satisfying the experimental bound~\cite{Baker2006} by 45 orders of magnitude.
No axion, Peccei--Quinn symmetry, or new particles are required.
This is the \emph{same exponential factor} as the vacuum energy suppression---two ``unrelated'' problems solved by one mechanism.

%% ============================================================
\section{Baryogenesis from BKT Freeze-Out}
\label{sec:baryogenesis}
%% ============================================================

\subsection{Sakharov conditions}

The BKT phase transition satisfies all three conditions:

\textbf{(1) B violation:}
Baryons are Hopf solitons with topological charge $Q \in \pi_3(S^2) = \mathbb{Z}$ (Paper~II).
Above $T_{\rm BKT}$, vortex-antivortex pairs unbind, topological protection fails, and soliton winding numbers can change: $\Delta B \neq 0$.
Below $T_{\rm BKT}$, vortex binding restores topological protection: $B$ conservation.

\textbf{(2) CP violation:}
The $A_2$ geometric CP phase $\delta_{\rm CP} = 68.3°$~\cite{Paper9} gives the Jarlskog invariant $J = 3.2\times10^{-5}$.

\textbf{(3) Departure from equilibrium:}
BKT critical slowing: $\tau_{\rm relax} = \tau_0\exp(b/\sqrt{|t|}) \to \infty$ as $t \to 0^+$.

\subsection{Freeze-out temperature}

Setting the color-mode vortex unwinding rate equal to the Hubble rate:
\begin{align}
\Gamma_{\rm unwind} &= \Gamma_0\,\exp\!\left(-\frac{2b_2}{\sqrt{t_f}}\right) = H(T_f)\,,\\
\Gamma_0 &= c/\varkappa = 10^{26}\,\text{s}^{-1}\,,\\
H(T \sim 150\,\text{MeV}) &\approx 4\times10^{-12}\,\text{s}^{-1}\,.
\end{align}
Solving:
\begin{equation}
t_{\rm freeze} = \left(\frac{\pi/3}{\ln(\Gamma_0/H)}\right)^2 = \left(\frac{1.047}{87.5}\right)^2 = 1.43\times10^{-4}\,.
\end{equation}

\textbf{Remarkably:} $t_{\rm freeze} \sim O(t_{\rm ESFT})$---the cosmological freeze-out occurs near the same BKT regime as all other ESFT predictions.

\subsection{Baryon asymmetry estimate}

\begin{equation}
\eta_B \sim J \times \frac{\Gamma_{\rm unwind}}{H}\bigg|_{t_f}\!\times\; W \times \frac{1}{g_*} \sim 3\times10^{-5}\times 1\times 0.01\times 0.06 \sim 10^{-8}\,,
\end{equation}
where $W \sim 0.01$ is the washout factor and $g_* = 17.25$ is the QCD-epoch relativistic degrees of freedom.

\textbf{Result:} $\eta_B^{\rm ESFT} \sim 10^{-8}$ vs.\ $\eta_B^{\rm obs} = 6.1\times10^{-10}$ (1.2 orders of magnitude).
For comparison, SM electroweak baryogenesis is $\sim 10$ orders of magnitude below observation.

\textbf{Unique prediction:} ESFT baryogenesis occurs at the QCD scale ($\sim 150\,$MeV), not the electroweak scale ($\sim 100\,$GeV).
This may be testable through baryon/antibaryon asymmetry signatures in heavy-ion collisions near the QGP phase boundary.

%% ============================================================
\section{Dark Matter Coherence Scale}
\label{sec:dm}
%% ============================================================

The BKT correlation length
\begin{equation}
\xi_{\rm BKT} = \varkappa\exp\!\left(\frac{\pi}{\sqrt{t}}\right) \sim 10\,\text{Mpc}
\end{equation}
explains the hierarchy of dark matter observations:
\begin{itemize}
\item $r \ll \xi$: BKT order complete $\to$ flat rotation curves.
\item $r \sim \xi$: BKT order decays $\to$ finite halo extent.
\item $r \gg \xi$: No coherence $\to$ cosmological homogeneity.
\end{itemize}

The background density $\rho_{\rm bg} = \sigma t\Lambda_{\rm core}/(\varkappa^2 c^2) = 4.1\,$GeV/cm$^3$ reproduces the Milky Way rotation curve ($v = 220\,$km/s, 0.1\%) and large SPARC galaxies ($R_d > 2\,$kpc) within $\pm 11\%$, all with zero free parameters.

\textbf{Testable:} Weak lensing surveys (Euclid, Rubin/LSST) should detect a sharp DM halo boundary at $\sim 10\,$Mpc, distinct from the NFW $1/r^3$ tail.

%% ============================================================
\section{BBN and Structure Formation}
\label{sec:bbn}
%% ============================================================

\subsection{BBN safety}

The ESFT modification to Newton's constant at the BBN epoch ($T \sim 1\,$MeV, $\rho \sim 10^6\,$kg/m$^3$) is:
\begin{equation}
\frac{\Delta G}{G}\bigg|_{\rm BBN} \sim \frac{\rho_{\rm BBN}}{\rho_{\rm crit}^{\rm ESFT}} \sim \frac{10^6}{10^{39}} \sim 10^{-33}\,,
\end{equation}
utterly negligible.
Standard BBN results (${}^4$He $\sim 25\%$, D/H $\sim 2.5\times10^{-5}$) are completely preserved.

\subsection{Structure formation}

ESFT predicts $G_{\rm eff}$ increasing at $z < 1$ as BKT locking strengthens, potentially explaining the $S_8$ tension (Planck $\sigma_8 = 0.83$ vs.\ lensing $\sigma_8 = 0.76$).
The transition scale is set by $\xi_{\rm BKT} \sim 10\,$Mpc.

%% ============================================================
\section{Cosmological Josephson Oscillation}
\label{sec:josephson}
%% ============================================================

The sine-Gordon master equation, when applied to the weak-coupling region between two BKT-ordered domains, reduces to the Josephson junction equation (Paper~IV, Sec.~\ref{sec:josephson}). At cosmological scales, the Josephson plasma frequency is BKT-renormalized:
\begin{equation}
\omega_J^{\rm cosmo} = \frac{\Lambda^2}{f}\exp\!\left(-\frac{b}{2\sqrt{t}}\right)\,.
\label{eq:omega_J_cosmo}
\end{equation}

For $b = \pi$ (gravitational relaxation mode):
\begin{equation}
\omega_J^{\rm cosmo} \approx 8\,{\rm MHz} \quad (\lambda \approx 37\,{\rm m})\,.
\end{equation}

\subsection{Signal characteristics}

The Josephson oscillation is a coherent, narrow-band, non-thermal process. The emitted radiation has:
\begin{itemize}
\item \textbf{Frequency:} $\sim 8\,$MHz ($b = \pi$) or $\sim 56\,$THz ($b = 2\pi/3$), depending on the dominant mode.
\item \textbf{Bandwidth:} $\Delta\nu/\nu \sim \sqrt{t} \approx 0.03$ (BKT fluctuations broaden the line).
\item \textbf{Spatial distribution:} Isotropic (domain walls are randomly oriented).
\item \textbf{Spectrum:} Non-thermal (monochromatic $\sin(\Delta\Theta)$ emission).
\end{itemize}

\subsection{Flux density estimate}

The energy density in Josephson oscillations at the domain wall is:
\begin{equation}
U_J \sim \frac{\Lambda^4}{2}\sin^2(\Delta\Theta_{\rm avg}) \times V_{\rm wall}/V_{\rm Hubble}\,.
\end{equation}
With domain walls at the BKT coherence scale $\xi_{\rm BKT} \sim 10\,$Mpc, the wall volume fraction is $\sim \koppa/\xi_{\rm BKT} \sim 10^{-44}$, and the average phase difference $\langle\sin^2\Delta\Theta\rangle \sim 1/2$ for random domains:
\begin{equation}
U_J \sim \frac{(73\,{\rm GeV})^4}{2} \times \frac{1}{2} \times 10^{-44} \sim 7\times10^{-19}\,{\rm GeV}^4\,.
\end{equation}

Converting to specific intensity at 8~MHz:
\begin{equation}
S_J \sim \frac{U_J c}{4\pi \Delta\nu} \sim 10^{-35}\,{\rm W\,m^{-2}\,Hz^{-1}} \sim 10^{-9}\,{\rm Jy}\,.
\end{equation}

This is $\sim 7$ orders of magnitude below LOFAR's confusion limit ($\sim 100\,\mu$Jy at 150~MHz). \textbf{The cosmological Josephson signal is undetectable with current technology.}

However, local overdensities (galaxy clusters, where $\rho \gg \bar{\rho}$) may enhance the signal through BKT stiffness amplification. A cluster at 100~Mpc with $\rho/\bar{\rho} \sim 100$ could boost the flux by $\sim 10^4$, to $\sim 10^{-5}\,$Jy---still below detection but approaching the sensitivity of next-generation instruments (SKA2, $\sim 10\,$nJy at 100~MHz).

\textbf{This is an honest null result for current observations, but a target for future radio astronomy.}

%% ============================================================
\section{The Bell Analogy}
\label{sec:bell}
%% ============================================================

The four-projection structure has a precise analogy to Bell's theorem~\cite{Bell1964}:

\begin{center}
\begin{tabular}{lll}
\toprule
& Bell's theorem & ESFT $A_2$ \\
\midrule
Hidden variable & $|\psi\rangle$ & $K_{A_2}$ \\
Measurement choice & Polarizer angle & Projection rule \\
Outputs & $E(a,b)$ & $\tau, \rho_\Lambda, T_{\rm CMB}, m_\nu$ \\
Classical bound & $|S| \leq 2$ & $\geq 4$ free parameters \\
Actual result & $|S| = 2\sqrt{2}$ & 0 free parameters, 96 OOM \\
\bottomrule
\end{tabular}
\end{center}

In Bell's theorem, a single quantum state produces correlations that no local hidden-variable model can reproduce.
In ESFT, a single Cartan matrix produces physical constants that no ``independent quantities'' model can reproduce---four quantities spanning 96 orders of magnitude, determined by zero adjustable parameters.

If all four quantities were truly independent, they would require at least four parameters.
The ESFT ``violation'' of this classical expectation---zero parameters for four quantities---is the strongest evidence for the unified framework.

Moreover, all four predictions move \emph{in lockstep} with $t$: increasing $t$ simultaneously shortens $\tau$, increases $\rho_\Lambda$, increases $T_{\rm CMB}$, and decreases $m_\nu$---with relative shifts rigidly fixed by the $b$ ratios.
This lock-step behavior is the analog of the $\cos 2\theta$ modulation in Bell correlations.

%% ============================================================
\section{Testable Predictions}
\label{sec:tests}
%% ============================================================

\begin{center}
\begin{tabular}{llll}
\toprule
Prediction & Value & Test & Timeline \\
\midrule
$\nu$ hierarchy & Normal & JUNO & $\sim$2027 \\
$m_{\nu_3}$ & 0.072\,eV & Cosmological & $\sim$2028 \\
DM boundary & $\sim$10\,Mpc & Euclid/Rubin & $\sim$2030 \\
$t$ spectroscopy & $1.14\times10^{-3}$ & Precision H spec. & $\sim$2030 \\
$\rho_\Lambda$ from $t$ & Factor $\sim 2$ & Cross-check & When $t$ measured \\
$S_8$ resolution & $G_{\rm eff}(z)$ & Lensing surveys & $\sim$2030 \\
$B$ at QCD scale & BKT signatures & RHIC/LHC-HI & Ongoing \\
\bottomrule
\end{tabular}
\end{center}

The strongest single test: precision spectroscopy determines $t$ independently; then $\rho_\Lambda$, $T_{\rm CMB}$, and $m_\nu$ are all predicted to $\sim 1\%$.
This constitutes a single measurement that simultaneously tests all four projections.

%% ============================================================
\section{Summary}
\label{sec:summary}
%% ============================================================

\begin{center}
\begin{tabular}{llllc}
\toprule
Quantity & $b_{\rm eff}$ & ESFT & Observed & Dev. \\
\midrule
$\tau/\tau_0$ & $\pi$ & $10^{41}$ & $10^{41}$ & $\checkmark$ \\
$\rho_\Lambda$ & $4\pi/3$ & $4\times10^{-48}$ & $3.6\times10^{-47}$ & $\times 9$ \\
$T_{\rm CMB}$ & $\pi/3$ & 2.62\,K & 2.725\,K & 3.9\% \\
$m_\nu$ & $\pi/6$ & 0.072\,eV & $\sim$0.050\,eV & 44\% \\
$H_0$ & --- & 72.5 & 67--73 & in band \\
$\theta_{\rm eff}$ & $4\pi/3$ & $10^{-55}$ & $<10^{-10}$ & $\checkmark$ \\
$\eta_B$ & BKT freeze & $10^{-8}$ & $6\times10^{-10}$ & $\times 30$ \\
\bottomrule
\end{tabular}
\end{center}

Seven cosmological/particle quantities, zero free parameters, one $2\times 2$ matrix.

%% ============================================================
\section{Conclusion}
\label{sec:conclusion}
%% ============================================================

We have shown that a single Cartan matrix $K_{A_2}$ with eigenvalues $\{1,3\}$, combined with the BKT reduced temperature $t = 1.1\times10^{-3}$, determines four fundamental physical quantities spanning 96 orders of magnitude through four different projection rules.
The vacuum energy (119 of 120 orders of magnitude explained), CMB temperature (3.9\%), neutrino mass ($\sim 0.07\,$eV), and dark matter relaxation time ($10^{41}\tau_0$) are not independent puzzles---they are four faces of the same $A_2$ BKT physics.

The apparent tension between $b = \pi$ and $b = 2\pi/3$ is resolved as two natural bases of the same matrix: Cartesian (field relaxation) and eigenbasis (vortex fugacity).

The same BKT factor simultaneously solves the strong CP problem ($\theta_{\rm eff} \sim 10^{-55}$), provides a baryogenesis mechanism ($\eta_B \sim 10^{-8}$, within 1.5 orders), and ensures BBN safety ($\Delta G/G \sim 10^{-33}$).

The strongest test: an independent spectroscopic determination of $t$ to 1\% precision simultaneously predicts $\rho_\Lambda$, $T_{\rm CMB}$, and $m_\nu$, each to $\sim$ factor-2 accuracy.
No other framework in physics connects atomic spectroscopy to cosmological constants through a single parameter.

\begin{acknowledgments}
This work was supported by no external funding.
AI-assisted tools were used for manuscript preparation; all physical reasoning and derivations are the author's.
\end{acknowledgments}

\begin{thebibliography}{25}

\bibitem{Paper1} J.~C.~P.~Won Ting, ``Topological soliton coherence scale in ESFT,'' submitted to Found.\ Phys.\ (2026).

\bibitem{Paper2} J.~C.~P.~Won Ting, ``Emergent gauge symmetry from Hopf soliton topology,'' submitted to Phys.\ Rev.\ D (2026).

\bibitem{Paper3} J.~C.~P.~Won Ting, ``Electron mass from topological string tension,'' submitted to Phys.\ Rev.\ Lett.\ (2026).

\bibitem{Paper7} J.~C.~P.~Won Ting, ``Fibonacci structure of fermion masses,'' in preparation (2026).

\bibitem{Paper9} J.~C.~P.~Won Ting, ``CKM and PMNS mixing from $A_2$ geometry,'' in preparation (2026).

\bibitem{Paper13} J.~C.~P.~Won Ting, ``Unified master equation of ESFT,'' in preparation (2026).

\bibitem{Fixsen2009} D.~J.~Fixsen, Astrophys.\ J.\ \textbf{707}, 916 (2009).

\bibitem{Planck2020} Planck Collaboration, Astron.\ Astrophys.\ \textbf{641}, A6 (2020).

\bibitem{Riess2022} A.~G.~Riess \textit{et al.}, Astrophys.\ J.\ Lett.\ \textbf{934}, L7 (2022).

\bibitem{Kosterlitz1974} J.~M.~Kosterlitz, J.\ Phys.\ C \textbf{7}, 1046 (1974).

\bibitem{Baker2006} C.~A.~Baker \textit{et al.}, Phys.\ Rev.\ Lett.\ \textbf{97}, 131801 (2006).

\bibitem{Bell1964} J.~S.~Bell, Physics \textbf{1}, 195 (1964).

\end{thebibliography}

\end{document}
